\clearpage
\SourceLeaf{105}%

\renewcommand{\BookRunningTitle}{Theorie der biquadratischen Reste}
\renewcommand{\Article}[1]{\par\medskip\Needspace{4\baselineskip}\pdfbookmark[1]{Artikel #1}{paper-five-art-#1}\centerline{\textbf{#1.}}\nobreak\smallskip}
\begin{center}
\textbf{THEORIE DER BIQUADRATISCHEN RESTE.}\par
\bigskip ZWEITE ABHANDLUNG.
\end{center}
\BookRule{.16\linewidth}
\Article{24}
In der ersten Abhandlung ist alles, was zur biquadratischen Klassifikation der Zahl $+2$ erforderlich ist, vollständig behandelt worden. Wir denken uns nämlich alle durch den Modul $p$ (der als Primzahl der Form $4n+1$ vorausgesetzt wird) nicht teilbaren Zahlen auf vier Sammlungen $A,B,C,D$ verteilt, je nachdem sie zur Potenz mit dem Exponenten $\frac14(p-1)$ erhoben modulo $p$ zu $+1$, $+f$, $-1$, $-f$ kongruent werden; dabei bezeichnet $f$ eine der beiden Wurzeln der Kongruenz $ff\equiv-1\ (\mathrm{mod.}\ p)$. Wir haben gefunden, dass die Entscheidung, welcher Sammlung die Zahl $+2$ zuzurechnen ist, von der Zerlegung der Zahl $p$ in zwei Quadrate abhängt: Setzt man $p=aa+bb$, wobei $aa$ das ungerade und $bb$ das gerade Quadrat bezeichnet, und nimmt man ferner an, dass die \emph{Vorzeichen} von $a,b$ so gewählt sind, dass $a\equiv1\ (\mathrm{mod.}\ 4)$, $b\equiv af\ (\mathrm{mod.}\ p)$, so muss die Zahl $+2$ zur Sammlung $A,B,C,D$ gehören, je nachdem $\frac12b$ die Form $4n$, $4n+1$, $4n+2$, $4n+3$ hat.

Daraus folgt auch unmittelbar eine Regel zur Klassifikation der Zahl $-2$. Da $-1$ zur Klasse $A$ gehört, wenn $\frac12b$ gerade ist, im ungeraden Fall aber zur Klasse $C$, gehört nach dem Satz in Artikel 7 die Zahl $-2$ zur Klasse $A,B,C,D$, je nachdem $\frac12b$ die Form $4n$, $4n+3$, $4n+2$, $4n+1$ hat.

Diese Sätze lassen sich auch folgendermaßen ausdrücken:

\SourceLeaf{106}%
\begin{center}
\begin{tabular}{c|c|c}
Es gehört&$+2$&$-2$\\\hline
zur Sammlung&\multicolumn{2}{c}{wenn $b$ modulo $8$ kongruent ist zu}\\\hline
$A$&$0$&$0$\\
$B$&$2a$&$6a$\\
$C$&$4a$&$4a$\\
$D$&$6a$&$2a$
\end{tabular}
\end{center}

Man sieht leicht, dass die so ausgesprochenen Sätze nicht mehr von der Bedingung $a\equiv1\ (\mathrm{mod.}\ 4)$ abhängen, sondern auch dann gelten, wenn $a\equiv3\ (\mathrm{mod.}\ 4)$, sofern die andere Bedingung $af\equiv b\ (\mathrm{mod.}\ p)$ erhalten bleibt.

Ebenso leicht sieht man, dass sich der Inhalt dieser Sätze elegant in einer einzigen Formel zusammenfassen lässt, nämlich:
\begin{quote}\itshape Werden $a$ und $b$ positiv genommen, so gilt stets\end{quote}
\[
b^{\frac12ab}\equiv a^{\frac12ab}\,2^{\frac14(p-1)}\ (\mathrm{mod.}\ p)
\]

\Article{25}
Sehen wir nun, wie weit die induktive Untersuchung auf die Klassifikation der Zahl $3$ hinweist. Setzt man die Tabelle in Artikel 11 weiter fort (wobei stets die kleinste primitive Wurzel genommen wird), so zeigt sie, dass $+3$ gehört
\begin{center}zur Sammlung\end{center}
\[
\begin{array}{r|r|r||r|r|r||r|r|r||r|r|r}
\multicolumn{3}{c||}{A\ \text{für}}&\multicolumn{3}{c||}{B\ \text{für}}&\multicolumn{3}{c||}{C\ \text{für}}&\multicolumn{3}{c}{D\ \text{für}}\\
p&a&b&p&a&b&p&a&b&p&a&b\\
13&-3&+2&17&+1&-4&37&+1&-6&5&+1&+2\\
109&-3&+10&29&+5&+2&61&+5&-6&41&+5&-4\\
181&+9&+10&53&-7&+2&73&-3&-8&149&-7&+10\\
193&-7&-12&89&+5&-8&97&+9&+4&173&+13&+2\\
229&-15&+2&101&+1&+10&157&-11&-6&&&\\
277&+9&+14&113&-7&-8&241&-15&-4&&&\\
&&&137&-11&-4&&&&&&\\
&&&197&+1&-14&&&&&&\\
&&&233&+13&+8&&&&&&\\
&&&257&+1&-16&&&&&&\\
&&&269&+13&+10&&&&&&\\
&&&281&+5&+16&&&&&&\\
&&&293&+17&+2&&&&&&
\end{array}
\]

\SourceLeaf{107}%
Zumindest auf den ersten Blick erkennen wir keinen einfachen Zusammenhang zwischen den Werten der Zahlen $a,b$, denen dieselbe Sammlung entspricht. Bedenken wir aber, dass sich eine entsprechende Entscheidung in der Theorie der quadratischen Reste für die Zahl $-3$ durch eine einfachere Regel treffen lässt als für die Zahl $+3$, so besteht die Hoffnung auf einen ebenso günstigen Erfolg in der Theorie der biquadratischen Reste. Wir finden, dass $-3$ zur Sammlung gehört
\[
\begin{array}{r|r|r||r|r|r||r|r|r||r|r|r}
\multicolumn{3}{c||}{A\ \text{für}}&\multicolumn{3}{c||}{B\ \text{für}}&\multicolumn{3}{c||}{C\ \text{für}}&\multicolumn{3}{c}{D\ \text{für}}\\
p&a&b&p&a&b&p&a&b&p&a&b\\
37&+1&-6&5&+1&+2&13&-3&+2&29&+5&+2\\
61&+5&-6&17&+1&-4&73&-3&-8&41&+5&-4\\
157&-11&-6&89&+5&-8&97&+9&+4&53&-7&+2\\
193&-7&-12&113&-7&-8&109&-3&+10&101&+1&+10\\
&&&137&-11&-4&181&+9&+10&197&+1&-14\\
&&&149&-7&+10&229&-15&+2&269&+13&+10\\
&&&173&+13&+2&241&-15&-4&293&+17&+2\\
&&&233&+13&+8&277&+9&+14&&&\\
&&&257&+1&-16&&&&&&\\
&&&281&+5&+16&&&&&&
\end{array}
\]
Hier zeigt sich die induktiv gefundene Regel unmittelbar. Die Zahl $-3$ gehört nämlich zur Sammlung
\[
\begin{array}{ll}
A,&\text{wenn }b\text{ durch }3\text{ teilbar ist, d.\,h.\ }b\equiv0\ (\mathrm{mod.}\ 3)\\
B,&\text{wenn }a+b\text{ durch }3\text{ teilbar ist, d.\,h.\ }b\equiv2a\ (\mathrm{mod.}\ 3)\\
C,&\text{wenn }a\text{ durch }3\text{ teilbar ist, d.\,h.\ }a\equiv0\ (\mathrm{mod.}\ 3)\\
D,&\text{wenn }a-b\text{ durch }3\text{ teilbar ist, d.\,h.\ }b\equiv a\ (\mathrm{mod.}\ 3)
\end{array}
\]

\Article{26}
Wir finden, dass die Zahl $+5$ der folgenden Sammlung zuzurechnen ist:
\[
\begin{array}{ll}
A\text{ für }p=&101,\ 109,\ 149,\ 181,\ 269\\
B\text{ für }p=&13,\ 17,\ 73,\ 97,\ 157,\ 193,\ 197,\ 233,\ 277,\ 293\\
C\text{ für }p=&29,\ 41,\ 61,\ 89,\ 229,\ 241,\ 281\\
D\text{ für }p=&37,\ 53,\ 113,\ 137,\ 173,\ 257
\end{array}
\]
Betrachten wir die Werte der Zahlen $a,b$, die den einzelnen $p$ entsprechen, so erkennen wir die Regel hier ebenso leicht wie bei der Klassifikation der Zahl $-3$. Wir gelangen nämlich zur Sammlung

\SourceLeaf{108}%
\[
\begin{array}{ll}
A,&\text{wenn }b\equiv0\ (\mathrm{mod.}\ 5)\\
B,&\text{wenn }b\equiv a\\
C,&\text{wenn }a\equiv0\\
D,&\text{wenn }b\equiv4a
\end{array}
\]

Offenbar erfassen diese Regeln alle Fälle, denn für $b\equiv2a$ oder $b\equiv3a\ (\mathrm{mod.}\ 5)$ wäre $aa+bb\equiv0$, Q. E. A., da nach Voraussetzung $p$ eine von $5$ verschiedene Primzahl ist.

\Article{27}
Ebenso weist die induktive Untersuchung, auf die Zahlen $-7$, $-11$, $+13$, $+17$, $-19$, $-23$ angewandt und hinreichend weit fortgesetzt, auf folgende Regeln hin:
\begin{center}Für die Zahl $-7$\end{center}
\[
\begin{array}{c|l}
A&a\equiv0,\quad\text{oder }b\equiv0\ (\mathrm{mod.}\ 7)\\
B&b\equiv4a,\ \text{oder }b\equiv5a\\
C&b\equiv a,\quad\text{oder }b\equiv6a\\
D&b\equiv2a,\ \text{oder }b\equiv3a
\end{array}
\]

\begin{center}Für die Zahl $-11$.\end{center}
\[
\begin{array}{c|l}
A&b\equiv0,\ 5a,\ \text{oder }6a\ (\mathrm{mod.}\ 11)\\
B&b\equiv a,\ 3a\text{ oder }4a\\
C&a\equiv0,\ \text{oder }b\equiv2a\text{ oder }9a\\
D&b\equiv7a,\ 8a\text{ oder }10a
\end{array}
\]

\begin{center}Für die Zahl $+13$.\end{center}
\[
\begin{array}{c|l}
A&b\equiv0,\ 4a,\ 9a\ (\mathrm{mod.}\ 13)\\
B&b\equiv6a,\ 11a,\ 12a\\
C&a\equiv0;\ b\equiv3a,\ 10a\\
D&b\equiv a,\ 2a,\ 7a
\end{array}
\]

\begin{center}Für die Zahl $+17$.\end{center}
\[
\begin{array}{c|l}
A&a\equiv0;\ b\equiv0,\ a,\ 16a\ (\mathrm{mod.}\ 17)\\
B&b\equiv2a,\ 6a,\ 8a,\ 14a\\
C&b\equiv5a,\ 7a,\ 10a,\ 12a\\
D&b\equiv3a,\ 9a,\ 11a,\ 15a
\end{array}
\]

\SourceLeaf{109}%
\begin{center}Für die Zahl $-19$.\end{center}
\[
\begin{array}{c|l}
A&b\equiv0,\ 2a,\ 5a,\ 14a,\ 17a\ (\mathrm{mod.}\ 19)\\
B&b\equiv3a,\ 7a,\ 11a,\ 13a,\ 18a\\
C&a\equiv0;\ b\equiv4a,\ 9a,\ 10a,\ 15a\\
D&b\equiv a,\ 6a,\ 8a,\ 12a,\ 16a
\end{array}
\]

\begin{center}Für die Zahl $-23$.\end{center}
\[
\begin{array}{c|l}
A&a\equiv0;\ b\equiv0,\ 7a,\ 10a,\ 13a,\ 16a\ (\mathrm{mod.}\ 23)\\
B&b\equiv2a,\ 3a,\ 4a,\ 11a,\ 15a,\ 17a\\
C&b\equiv a,\ 5a,\ 9a,\ 14a,\ 18a,\ 22a\\
D&b\equiv6a,\ 8a,\ 12a,\ 19a,\ 20a,\ 21a
\end{array}
\]

\Article{28}
Die so durch induktive Untersuchung gewonnenen besonderen Sätze bestätigen sich, soweit diese Untersuchung fortgesetzt wird, und zeigen eine besonders schöne Gestalt der Kriterien. Vergleicht man sie aber miteinander, um allgemeine Schlüsse daraus zu ziehen, so bieten sich sogleich die folgenden Beobachtungen dar.

Die Kriterien dafür, welcher Sammlung die Primzahl $\pm q$ zuzurechnen ist (mit dem oberen oder unteren Vorzeichen, je nachdem $q$ die Form $4n+1$ oder $4n+3$ hat), hängen von den Formen der Zahlen $a,b$ ab, wenn man diese bezüglich des Moduls $q$ miteinander vergleicht. Es gilt nämlich:

\Needspace{11\baselineskip}
I. Ist $a\equiv0\ (\mathrm{mod.}\ q)$, so gehört $\pm q$ zu einer bestimmten Sammlung, und zwar zu $A$ für $q=7,\ 17,\ 23$ und zu $C$ für $q=3,\ 11,\ 13,\ 19$. Daraus entsteht die Vermutung, dass der erste Fall allgemein gilt, wenn $q$ die Form $8n\pm1$ hat, der zweite dagegen, wenn $q$ die Form $8n\pm3$ hat. Übrigens sind die Sammlungen $B$ und $D$ bereits ohne induktive Untersuchung ausgeschlossen, wenn der Wert von $a$ durch $q$ teilbar ist. Dann gilt $p\equiv bb\ (\mathrm{mod.}\ q)$, das heißt, $p$ ist ein quadratischer Rest von $q$, und daher muss nach dem Fundamentalsatz $\pm q$ ein quadratischer Rest von $p$ sein.

II. Ist dagegen $a$ nicht durch $q$ teilbar, so hängt das Kriterium vom Wert des Ausdrucks $\frac ba\ (\mathrm{mod.}\ q)$ ab. Dieser Ausdruck lässt zwar $q$ verschiedene Werte zu, nämlich $0,\ 1,\ 2,\ 3\dotsfour q-1$; hat aber $q$ die Form $4n+1$, so müssen die beiden
\SourceLeaf{110}%
Werte des Ausdrucks $\surd-1\ (\mathrm{mod.}\ q)$ ausgeschlossen werden. Diese können offenbar keine Werte des Ausdrucks $\frac ba\ (\mathrm{mod.}\ q)$ sein, da $p=aa+bb$ stets als eine von $q$ verschiedene Primzahl vorausgesetzt wird. Somit ist die Anzahl der zulässigen Werte des Ausdrucks $\frac ba\ (\mathrm{mod.}\ q)$ $=q-2$ für $q\equiv1\ (\mathrm{mod.}\ 4)$, während sie $=q$ für $q\equiv3\ (\mathrm{mod.}\ 4)$ bleibt.

Diese Werte verteilen sich nun auf vier Klassen: Einige, die wir unbestimmt mit $\alpha$ bezeichnen, entsprechen der Sammlung $A$, andere, mit $\beta$ bezeichnet, der Sammlung $B$, andere, $\gamma$, der Sammlung $C$, und schließlich die übrigen, $\delta$, der Sammlung $D$. Die Zahl $\pm q$ ist also der Sammlung $A,B,C,D$ zuzurechnen, je nachdem $b\equiv\alpha a$, $b\equiv\beta a$, $b\equiv\gamma a$, $b\equiv\delta a\ (\mathrm{mod.}\ q)$ gilt.

Doch das \emph{Gesetz} dieser Verteilung scheint tiefer verborgen, obwohl einige allgemeine Eigenschaften leicht zu erkennen sind. Die Anzahl der Werte ist in drei Klassen gleich, nämlich $=\frac14(q-1)$ oder $\frac14(q+1)$, während sie in einer (und zwar in derselben, die der Sammlung mit dem Kriterium $a\equiv0$ entspricht) um eins kleiner ist. Dadurch wird die Gesamtzahl der verschiedenen Kriterien für jede einzelne Sammlung gleich, nämlich $=\frac14(q-1)$ oder $\frac14(q+1)$. Ferner bemerken wir, dass $0$ stets in der ersten Klasse (unter den $\alpha$) vorkommt und dass die Ergänzungen der Zahlen $\alpha,\beta,\gamma,\delta$ zu $q$, also $q-\alpha$, $q-\beta$, $q-\gamma$, $q-\delta$, jeweils in der ersten, vierten, dritten und zweiten Klasse liegen. Schließlich sehen wir, dass die Werte der Ausdrücke $\frac1\alpha$, $\frac1\beta$, $\frac1\gamma$, $\frac1\delta\ (\mathrm{mod.}\ q)$ zur ersten, vierten, dritten und zweiten Klasse gehören, wenn das Kriterium $a\equiv0$ der Sammlung $A$ entspricht, dagegen jeweils zur dritten, zweiten, ersten und vierten Klasse, wenn sich das Kriterium $a\equiv0$ auf die Sammlung $C$ bezieht. Doch darauf beschränkt sich ungefähr, was sich durch induktive Untersuchung erreichen lässt, sofern wir uns nicht allzu kühn anmaßen wollen, das vorwegzunehmen, was später aus den eigentlichen Quellen geschöpft werden wird.
